\section{Problem Setup}
\label{sec:problem}

We now state the detection and grading problem formally, fixing the notation and
evaluation protocol used throughout our experiments (Sec.~\ref{sec:application}).
Figure~\ref{fig:pipeline} summarizes the pipeline.

\subsection{Flat multi-class formulation}
\label{sec:problem:flat}

\noindent\textbf{Clips and labels.}
A recording session provides temporally synchronized video and EEG. Around each
expert-annotated event we extract a clip $c=(\mathbf{V},\mathbf{E})$, where
$\mathbf{V}\in\mathbb{R}^{T\times H\times W\times 3}$ is a short RGB video segment
and $\mathbf{E}\in\mathbb{R}^{C\times L}$ the corresponding $C$-channel EEG trace.
Matched non-seizure clips are sampled from seizure-free intervals of the same
sessions. Each clip carries a severity label $y\in\{0,1,2,3,4\}$: $0$
denotes non-seizure and $1$--$4$ the ictal stages (Racine~\cite{racine1972}
Stage\,2--5). Racine Stage\,1 (mouth and facial automatisms) is annotated in this
cohort but occurs in only six clips, too few to train or evaluate a class on, so we
exclude it. The modeled label space begins at Stage\,2 as a consequence of the
data, and not because the scale starts there. The scale itself was derived from electrically
kindled seizures, and its transfer to chemically induced models is not
automatic~\cite{vanerum2019racine}; we use the grades our annotator assigned and do not re-derive a
model-specific scale. Each clip also belongs to a subject $s\in\mathcal{S}$ and a session
$r\in\mathcal{R}$. The corpus pools clips from all subjects.
Figure~\ref{fig:setup} shows a representative frame of the recording setup;
because camera geometry differs per cage, each clip is cropped per cage before
subjects are pooled.

% fig:setup now lives in sec/1_intro.tex so it prints on p.2 alongside its first
% mention in the contributions list.

\noindent\textbf{Task granularities.}
We study three classification tasks of increasing difficulty, each obtained by
collapsing the stage label through a map $g_k:\{0,\dots,4\}\to\{0,\dots,k-1\}$:
\begin{itemize}[leftmargin=1.2em,itemsep=1pt,topsep=2pt]
  \item \textbf{Detection} ($k{=}2$): $g_2(y)=\mathbf{1}[y>0]$: seizure \vs\
        non-seizure.
  \item \textbf{Coarse grading} ($k{=}3$): non-seizure / mild (Stage\,2--3) /
        severe (Stage\,4--5).
  \item \textbf{Fine grading} ($k{=}5$): the full non-seizure / Stage\,2--5 scale.
\end{itemize}

\noindent\textbf{Modalities.}
Each task pairs two \emph{unimodal} predictors that we compare head-to-head:
a video model $f_{\mathrm{V}}:\mathbf{V}\mapsto\hat{y}$ and an EEG model
$f_{\mathrm{E}}:\mathbf{E}\mapsto\hat{y}$. The video model consumes a fixed number
of sampled frames per clip. The EEG model is trained on short fixed windows and
forms a clip decision by mean-pooling the window posteriors. We keep the two
modalities \emph{separate} by design: our goal is to measure how far each single
modality gets on its own, so a fused predictor $f_{\mathrm{V,E}}$ is out of scope
and left to future work (Sec.~\ref{sec:open}).

\noindent\textbf{Evaluation protocol.}
Seizure burden is highly uneven across subjects, so a single pooled score can
mislead. We therefore adopt three safeguards throughout. (i)~\emph{Session-disjoint
splits}: all clips from a session $r$ fall entirely in train or validation, so no
clip leaks across the split. The split seed is shared between the
video and EEG models so their numbers are directly comparable. (ii)~\emph{Class
imbalance} is handled with inverse-frequency-weighted cross-entropy in place of
resampling. (iii)~We report \emph{macro-averaged} metrics (macro-F1 and, for
detection, AUROC) together with the Matthews correlation coefficient (MCC);
App.~\ref{sec:metrics} states each of them formally, along with the
fold- and seed-aggregation convention.
The two summarize different parts of the same confusion matrix: macro-F1 weights
every stage equally, so the rare severe classes dominate it, while MCC is computed
from the whole table and so follows the frequent classes. We report both so
that no conclusion rests on one weighting, and always pair the pooled score with a
per-subject breakdown, since aggregate accuracy is dominated by the few high-burden animals.
Sec.~\ref{sec:application} additionally reports a stricter protocol:
held-out-\emph{subject} generalization via subject-disjoint cross-validation.

\subsection{Conditional ordinal severity grading}
\label{sec:problem:ordinal}

The $k{=}5$ task above places non-seizure and the four ictal stages on a single
flat axis, which misstates the label space in two ways. First, ``non-seizure'' is
not a point on the Racine scale~\cite{racine1972}: the scale grades the motor
manifestations of a seizure, so even its lowest rung is an ictal event, and the
absence of an event is a different kind of object altogether. Second, a flat
softmax treats the four stages as exchangeable, so predicting Stage\,5 for a
Stage\,2 clip costs exactly what predicting Stage\,3 costs, even though severity is
ordered and the two errors differ in clinical consequence.

We therefore also study severity as a \emph{conditional ordinal} problem. Detection
and grading factorize,
\begin{equation}
p(y \mid c) \;=\; \underbrace{p(\text{seizure} \mid c)}_{\text{detection}}
\;\cdot\;
\underbrace{p(s \mid c,\ \text{seizure})}_{\text{severity grading}},
\label{eq:conditional}
\end{equation}
and we condition on the first factor: grading is posed on seizure clips alone, over
the ordered stage set $s\in\{\mathrm{S2}\prec\mathrm{S3}\prec\mathrm{S4}\prec
\mathrm{S5}\}$. This matches how the annotation is produced, since an expert
first confirms an event and then grades it, and it removes the 12{,}197 non-seizure clips
that dominate the flat 5-class objective.

\noindent\textbf{Ordinal head.}
Instead of four independent logits, an ordinal model predicts the $K{-}1$ cumulative
decisions $p(s \succ \mathrm{S2})$, $p(s \succ \mathrm{S3})$, $p(s \succ
\mathrm{S4})$, whose monotone structure encodes the ordering directly; the stage
follows from how many thresholds are exceeded. The nominal four-way classifier
trained on the same clips is the control, so any difference is attributable to the
ordinal structure alone.

\noindent\textbf{Metrics.}
Macro-F1 keeps continuity with the flat tasks, but it is order-blind, so we pair it
with measures that are not: the mean absolute stage error
$\frac{1}{N}\sum_i|\hat{s}_i-s_i|$, quadratic weighted $\kappa$, and the rate of
\emph{large-stage errors} $|\hat{s}-s|\ge 2$, which is the quantity a screening tool
most needs to keep small. Because the deployment question is a new animal, we
evaluate under the subject-disjoint five-fold protocol of
Sec.~\ref{sec:crossmodal}.
